
\section{Cronograma}

De acordo com a metodologia proposta, o cronograma apresentado na Tabela~\ref{tab:cronograma} organiza as etapas previstas para os próximos cinco meses, distribuindo as atividades de forma sequencial e alinhada aos objetivos do trabalho.

\begin{table*}[!ht]
\centering
\small % reduz o tamanho da fonte para caber na página
\caption{Cronograma de atividades ao longo de cinco meses.}
\label{tab:cronograma}
\renewcommand{\arraystretch}{1.2}
\begin{tabular}{|p{6.5cm}|c|c|c|c|c|}
\hline
\textbf{Atividade} & \textbf{Fevereiro} & \textbf{Março} & \textbf{Abril} & \textbf{Maio}  & \textbf{Junho} \\ \hline

Revisão dos algoritmos e requisitos             & x & x &   &   &   \\ \hline
Implementação dos algoritmos simples (Python)   & x & x &   &   &   \\ \hline
Comparação visual e ajustes                     & x & x &   &   &   \\ \hline
Configuração ROS2, Gazebo e PX4 SITL            & x & x &   &   &   \\ \hline
Cenário 1 (livre)                               &   &   & x &   &   \\ \hline
Cenário 2 (obstáculos estáticos)                &   &   & x &   &   \\ \hline
Cenário 3 (obstáculos dinâmicos)                &   &   & x &   &   \\ \hline
Execuções (10--20 simulações por cenário)       &   &   & x & x &   \\ \hline
Validação dos dados e reexecuções               &   &   &   & x &   \\ \hline
Cálculo das métricas e gráficos                 &   &   &   & x &   \\ \hline
Análise comparativa e conclusões                &   &   &   & x & x \\ \hline
Redação (resultados, conclusão, discussão)      &   &   &   & x & x \\ \hline
Revisão final e preparação da defesa            &   &   &   &   & x \\ \hline

\end{tabular}
\end{table*}


% \begin{declarations}

% \begin{acknowledgements}
% ESTA DECLARAÇÃO É OPCIONAL. Este é um texto de agradecimentos com várias linhas. Lorem ipsum dolor sit amet, consectetur adipiscing elit, sed do eiusmod tempor incididunt ut labore et dolore magna aliqua. Ut enim ad minim veniam, quis nostrud exercitation ullamco laboris nisi ut aliquip ex ea commodo consequat.
% \end{acknowledgements}


% \begin{funding}
% ESTA DECLARAÇÃO É OPCIONAL. Esta pesquisa foi financiada por lorem ipsum dolor sit amet, consectetur adipiscing elit.
% \end{funding}

% \begin{contributions}
% ESTA DECLARAÇÃO É OBRIGATÓRIA. 
% \textit{Considerar o Código de Conduta da SBC.
% Participação em autoria: espera-se que todos os autores de um trabalho publicado 
% ou submetido para publicação, tenham tido efetiva participação no respectivo trabalho.}
% Sugerimos que os autores descrevam sua contribuição usando 
% a Taxonomia CRediT (\href{https://credit.niso.org/}{https://credit.niso.org/}) como neste exemplo: 
% JV contribuiu para a concepção deste estudo. CB, RP e CM realizaram os experimentos. 
% JV é o principal contribuidor e escritor deste manuscrito. 
% Todos os autores leram e aprovaram o manuscrito final. 
% \end{contributions}

% \begin{interests}
% ESTA DECLARAÇÃO É OBRIGATÓRIA. 
% Se não houver conflitos de interesse, os autores devem declarar: 
% ``Os autores declaram que não têm nenhum conflito de interesses''. 
% Caso contrário, a declaração deve ser: 
% ``Os autores declaram que têm os seguintes conflitos de interesses: 
% lorem ipsum dolor sit amet, consectetur adipiscing elit.''
% \end{interests}

% \begin{materials}
% ESTA DECLARAÇÃO É OBRIGATÓRIA. 
% \textit{Considerar o Código de Conduta da SBC. 
% Reprodutibilidade de resultados de pesquisa: 
% nos casos pertinentes, recomenda-se que artigos indiquem 
% a disponibilidade pública de material utilizado nas revisões, 
% de modo a facilitar a reprodução dos respectivos resultados por outros pesquisadores.}
% %
% Se os autores disponibilizaram seus dados e/ou códigos abertamente, a declaração deve ser: 
% ``Os dados e outros materiais criados e/ou usados nesta revisão da literatura
% estão disponíveis em \ldots''. 
% Caso contrário, a declaração deve ser: 
% ``Os dados e outros materiais criados e/ou usados nesta revisão da literatura 
% serão disponibilizados mediante solicitação''. 
% %
% Todo o material suplementar disponibilizado publicamente pelos autores aparecerá 
% na página de publicação do artigo em caso de aceitação. 
% \end{materials}

% \begin{aitools}
% ESTA DECLARAÇÃO É OBRIGATÓRIA.  
% \textit{Considerar o Código de Conduta da SBC.
% Uso de Inteligência Artificial (IA) Generativa: 
% a utilização de ferramentas e tecnologias de IA Generativa para geração de conteúdos, 
% na escrita e/ou revisão do conteúdo de artigos, deve ser declarada explicitamente no trabalho.} 
% A declaração deve ocorrer nesta seção, definida especificamente para este fim.
% Deve-se listar as ferramentas (e suas versões) e descrever onde foram empregadas, por exemplo, 
% textos, tabelas, gráficos, citações, etc. 
% Essas ferramentas \textbf{não} podem ser listadas como autores de um artigo. 
% O uso de tais ferramentas não exime os autores da responsabilidade 
% sobre todo o seu conteúdo, inclusive no caso de ser identificado plágio.
% %
% Se os autores não utilizaram tais ferramentas, a declaração deve ser: 
% ``Os autores não utilizaram ferramentas de IA generativa no desenvolvimento do estudo''.
% %
% \end{aitools}

% \begin{furtherinformation}
% Manuscritos devem indicar explicitamente como as questões éticas da pesquisa foram abordadas e gerenciadas, 
% incluindo a aprovação da pesquisa por um comitê de ética, quando aplicável.
% \end{furtherinformation}

% \end{declarations}

